The error-structure and blink-type-recall findings indicate that the proposed robust epoch-aware threshold produced a clean but conservative detection profile. In Table~\ref{tab:error-structure}, BLINKER-concat was strongly false-positive-heavy, indicating over-detection, whereas MNE-annot and both proposed configurations were false-negative-heavy; therefore, the proposed methods traded some recall for far fewer false positives. This pattern is important because it suggests that the robust epoch-aware threshold is preferable in applications where false positives are costly, although it may miss a subset of true blink events. The blink-type analysis in Table~\ref{tab:blink_type_recall} further showed that recall was high for normal blinks and slightly lower for long blinks for three of the four methods, with Proposed-Med achieving the best recall among the conservative methods for both normal blinks (0.8548) and long blinks (0.8230). Together, these results imply that Proposed-Med offers a balanced conservative alternative with reduced over-detection, while long blinks remain a small residual weakness that merits future attention.
